%% sections/introduction.tex — v12: compressed to ~150 words

\section{Introduction}

Therapeutic circRNA offers advantages over linear mRNA in stability and tunable immunogenicity \citep{wesselhoeft2018,chen2019}, yet no computational platform evaluates circRNA-specific dimensions: modification effects on half-life ($\Psi$: ${\sim}15$\,h vs unmodified: 6.24\,h), circRNA immunogenicity pathways, IRES-mediated translation, or miRNA sponge activity. Existing tools address individual modalities---NetMHCpan for MHC binding \citep{reynisson2020}, ADMETlab for ADMET screening \citep{xiong2024}, Monolix for population PK \citep{dumont2014}---but none operates across both small-molecule and circRNA modalities. We developed Confluencia to integrate these evaluations in a single platform, producing unified Go/Conditional/No-Go decisions that guide our wet-lab prioritization. Three design choices address the data-limited setting (Fig~\ref{fig:overview}): (1) RNACTM, a literature-parameterized compartment simulation for circRNA dynamics; (2) sample-size-adaptive Mixture-of-Experts (MOE) \citep{jacobs1991} for binding prediction; (3) uncertainty-adaptive 5D evaluation that auto-downweights under-constrained dimensions via $(1{-}u)^2$.